\documentclass{article}
\usepackage{arxiv}

\usepackage[utf8]{inputenc}
\usepackage[T1]{fontenc}
\usepackage{hyperref}
\usepackage{url}
\usepackage{booktabs}
\usepackage{amsfonts}
\usepackage{amssymb}
\usepackage{amsmath}
\usepackage{amsthm}
\usepackage{mathtools}
\usepackage{bm}
\usepackage{microtype}
\usepackage{graphicx}
\usepackage{xcolor}
\usepackage{enumitem}
\usepackage{natbib}

% --------------------------------------------------------------------
%  Macros
% --------------------------------------------------------------------
\newcommand{\R}{\mathbb{R}}
\DeclareMathOperator{\SE}{SE}
\DeclareMathOperator{\Sim}{Sim}
\DeclareMathOperator{\SO}{SO}
\DeclareMathOperator{\Ad}{Ad}
\DeclareMathOperator{\img}{im}
\DeclareMathOperator{\rank}{rank}
\DeclareMathOperator{\spann}{span}
\newcommand{\norm}[1]{\left\lVert #1 \right\rVert}
\newcommand{\inner}[2]{\left\langle #1,\, #2 \right\rangle}
\newcommand{\sheaf}{\mathcal{F}}
\newcommand{\Cochain}{C}
\newcommand{\Lap}{\mathbf{L}}
\newcommand{\coboundary}{\delta}
\newcommand{\restr}[2]{\sheaf_{#1 \,\trianglelefteq\, #2}}
\newcommand{\Sections}{\Gamma}
\newcommand{\Hcoh}{H}
\newcommand{\Mdyn}{M_{\mathrm{dyn}}}
\newcommand{\Lcons}{\mathcal{L}_{\mathrm{cons}}}
\newcommand{\Ldyn}{\mathcal{L}_{\mathrm{dyn}}}
\newcommand{\Lstatic}{\mathcal{L}_{\mathrm{static}}}
\newcommand{\Lanchor}{\mathcal{L}_{\mathrm{anchor}}}
\newcommand{\todo}[1]{\textcolor{red}{[TODO: #1]}}

\theoremstyle{definition}
\newtheorem{definition}{Definition}
\theoremstyle{plain}
\newtheorem{theorem}{Theorem}
\newtheorem{proposition}{Proposition}
\newtheorem{lemma}{Lemma}
\theoremstyle{remark}
\newtheorem{remark}{Remark}

% --------------------------------------------------------------------
\title{Multi-Frame Consistency for Dynamic Pointmap Regression:\\
A Sheaf-Theoretic Training Objective}
% Working title -- alternatives:
%   "Baking Multi-Frame Consistency into Dynamic Pointmap Regression"
%   "Lifting Pairwise Pointmaps to Multi-Frame Consistency via Sheaf Cohomology"

\author{%
  Author One \\
  Department / Institution \\
  \texttt{author1@institution.edu} \\
  \And
  Author Two \\
  Department / Institution \\
  \texttt{author2@institution.edu} \\
}

\renewcommand{\headeright}{}
\renewcommand{\undertitle}{}

\begin{document}
\maketitle

\begin{abstract}
Feed-forward pointmap regressors in the DUSt3R family map an image
\emph{pair} to per-pixel 3D coordinates in a shared frame, recovering
depth, relative pose, and dense correspondence at once.
$\mathrm{D}^2$USt3R extends this to dynamic scenes with static--dynamic
aligned pointmaps, but, like its predecessors, it is fundamentally a
\emph{two-frame} model: reconstructing $N>2$ views relies on a separate,
post-hoc global-alignment optimization, and the authors note that
extending the model to video would require an additional global step
that the network itself does not provide. We argue this multi-frame
consistency should be a \emph{training} signal, not an inference-time
patch. The obstacle is that naive multi-frame consistency penalizes
\emph{all} cross-view disagreement, which forces moving objects into a
single rigid placement and smears their geometry---precisely the failure
mode that motivated $\mathrm{D}^2$USt3R. We resolve this by formulating
multi-frame consistency on a \emph{cellular sheaf} over the view graph,
whose Hodge decomposition splits cross-view disagreement into a
pose-explainable (coboundary) part that a camera correction can absorb
and an irreducible (harmonic) part that only object motion produces. Our
objective drives the coboundary part to zero---enforcing static
consistency across the whole tuple---while \emph{permitting} the harmonic
part on dynamic regions, lifting $\mathrm{D}^2$USt3R's per-pair
static/dynamic loss split to the multi-frame view graph. A static-scene
anchor, on which harmonic mass must vanish, together with an explicit
scale-gauge constraint, keeps the objective non-degenerate. The method is
a training-time term that sits \emph{on top of} $\mathrm{D}^2$USt3R rather
than competing with it.
\todo{Fill in headline numbers once experiments are run.}
\end{abstract}

% ====================================================================
\section{Introduction}
% ====================================================================

Pointmap regression has reshaped multi-view 3D reconstruction. In place
of the classical pipeline---detect correspondences, estimate poses,
triangulate, fuse---models in the DUSt3R
family~\citep{wang2024dust3r,leroy2024grounding} learn a single
feed-forward map from an image pair $(I^1,I^2)$ to two dense
\emph{pointmaps}, per-pixel 3D coordinates expressed in a common camera
frame, so that depth, relative pose, and dense correspondence are read off
a single prediction. For dynamic scenes, where object motion breaks the
rigid-alignment assumption, MonST3R~\citep{zhang2025monst3r} retrains on
moving content and $\mathrm{D}^2$USt3R~\citep{han2026enhancing} introduces
static--dynamic aligned pointmaps that warp static pixels by camera pose
and dynamic pixels by optical flow, markedly sharpening reconstruction in
dynamic regions. We build on $\mathrm{D}^2$USt3R most directly.

\paragraph{The gap: consistency is two-frame.}
These models are inherently \emph{pairwise}. Recovering a scene from
$N>2$ views requires a separate \emph{fusion} stage that reconciles the
overlapping pairwise predictions into one globally consistent geometry,
handled today by a per-scene global-alignment optimization over per-view
depths, poses, and intrinsics~\citep{wang2024dust3r,zhang2025monst3r}, or
by a learned recurrent state~\citep{wang2025continuous,wang20253d}. In
either case, multi-frame consistency is supplied \emph{after} the
network, by a generic optimizer that the per-pair predictions were never
trained to make easy. $\mathrm{D}^2$USt3R is explicit that its model is
non-generative, takes two frames, and that extending it to video ``would
require an additional global optimization technique.'' That technique is
exactly what is missing, and we contend it should be folded into training:
a network that emits pairwise predictions which are already
multi-frame-consistent removes the need to repair them downstream.

\paragraph{Why naive multi-frame consistency is wrong for dynamics.}
The obvious move---add a loss penalizing cross-view pointmap disagreement
over $N>2$ frames---fails in exactly the regime we care about. A moving
object photographed at two times has \emph{no} single rigid placement
consistent with both views; penalizing its disagreement forces the
network to pick one, smearing the object across its trajectory. This is
the MonST3R failure $\mathrm{D}^2$USt3R diagnoses: aligning all pointmaps
by a single global rigid transform corrupts moving regions. A usable
multi-frame consistency objective must therefore \emph{distinguish} the
disagreement a camera correction can explain (which should vanish) from
the disagreement only motion can produce (which must be allowed). That
distinction is not a heuristic threshold; it is an exact algebraic split.

\paragraph{Our viewpoint.}
Index the $N$ views by the vertices of a \emph{view graph} $G=(V,E)$, with
an edge per related pair. The per-pair predictions assign a local geometry
to each view and a relative transform to each edge---the data of a
\emph{cellular sheaf}~\citep{curry2014sheaves,ghrist2014elementary} whose
stalks are per-view geometries and whose restriction maps are the
predicted transforms. A globally consistent reconstruction is a
\emph{global section} of this sheaf: an assignment that agrees, edge by
edge, after restriction. The sheaf Laplacian's Hodge decomposition then
splits any cross-view residual into a \emph{coboundary} part---removable
by a per-view (pose/frame) correction---and a \emph{harmonic} part, the
first cohomology $\Hcoh^1(G;\sheaf)$, that no per-view correction can
remove~\citep{hansen2019toward,robinson2017sheaves}. Camera-pose error and
drift live entirely in the coboundary; irreducible cross-view
disagreement, the signature of object motion, lives in the harmonic part.
This is precisely the split a dynamics-aware consistency loss needs, and
it generalizes $\mathrm{D}^2$USt3R's per-pair static/dynamic supervision
to the whole tuple.

\paragraph{From detector to objective.}
We stress what this paper is \emph{not}. A natural temptation is to use
the harmonic part as an unsupervised, flow-free \emph{detector} of moving
regions. We do not pursue that: optical-flow motion segmentation, given an
accurate backbone, is already strong, and a detector competes with it on
its own ground. Our contribution is instead a \emph{training objective}
that uses the coboundary/harmonic split to shape the network's
predictions, so that static structure becomes multi-frame-consistent by
construction while dynamic structure is left free. The role is
complementary to $\mathrm{D}^2$USt3R, not competing.

\paragraph{Non-degeneracy.}
A consistency objective is gameable: the network can drive cross-view
disagreement to zero by collapsing scale, over-smoothing geometry, or
agreeing trivially. We close this with two constraints. First, a
\emph{static-scene anchor}: on sequences known to be rigid, the harmonic
part must vanish everywhere, which forbids the network from explaining
static structure as motion. Second, an explicit \emph{scale-gauge}
constraint pinning the one global degree of freedom (the uniform log-scale
direction) that no differential consistency term can determine and that
otherwise admits a collapse to a point.

\paragraph{Contributions.}
\begin{enumerate}[leftmargin=1.4em,itemsep=0.15em,topsep=0.2em]
  \item We recast multi-frame pointmap consistency as the search for a
        global section of a cellular sheaf on the view graph, and identify
        its coboundary/harmonic split with the static/dynamic split that a
        dynamics-aware consistency loss requires
        (Section~\ref{sec:method}). This lifts $\mathrm{D}^2$USt3R's
        per-pair static/dynamic supervision to $N>2$ frames.
  \item We turn this split into a training objective that enforces
        multi-frame static consistency while permitting dynamic motion,
        regularized by a static-scene anchor and a scale gauge that make
        it non-degenerate, and that attaches to a frozen-encoder
        $\mathrm{D}^2$USt3R backbone as an added term
        (Section~\ref{sec:method}).
  \item We design an experimental protocol (Section~\ref{sec:experiments})
        that tests whether the objective improves multi-frame
        reconstruction over the two-frame baseline, isolates the
        coboundary/harmonic split against naive consistency, and verifies
        non-degeneracy on static-scene controls.
\end{enumerate}

% ====================================================================
\section{Preliminaries}
\label{sec:prelim}
% ====================================================================

We fix notation for pointmaps and multi-frame fusion
(Section~\ref{sec:prelim-pointmaps}), recall cellular sheaves and their
Hodge/Laplacian structure (Section~\ref{sec:prelim-sheaves}), and assemble
the reconstruction sheaf the objective is built on
(Section~\ref{sec:recon-sheaf}).

% --------------------------------------------------------------------
\subsection{Pointmaps and multi-frame fusion}
\label{sec:prelim-pointmaps}
% --------------------------------------------------------------------

Let $I^1,\dots,I^N$ be images of resolution $W\times H$. A \emph{pointmap}
$X^{n,m}\in\R^{W\times H\times 3}$ assigns to each pixel of image $n$ its
3D location in the camera frame of image $m$. For a pair $(n,m)$ the
backbone~\citep{wang2024dust3r,zhang2025monst3r,han2026enhancing} regresses
$X^{n,m}$ and $X^{m,m}$ with per-pixel confidences $C^{n,m}$, from which
depth, the relative camera transform, and dense correspondence follow.
Writing $\pi_K^{-1}(p,d)=d\,K^{-1}\tilde p$ for back-projection of pixel
$p$ at depth $d$ through intrinsics $K$, each view $n$ with depth $D_n$,
pose $g_n\in\SE(3)$, and intrinsics $K_n$ contributes world-frame points
\begin{equation}
  \chi_n(p) \;=\; g_n \cdot \pi_{K_n}^{-1}\!\big(p,\,D_n(p)\big).
  \label{eq:chi}
\end{equation}

\paragraph{$\mathrm{D}^2$USt3R supervision (what we extend).}
$\mathrm{D}^2$USt3R supervises a pair with a dynamic-mask-gated split: a
static regression term aligning static pixels by camera pose, and a
dynamic alignment term aligning moving pixels via optical-flow
correspondence, gated by a dynamic mask $\Mdyn$ derived from the
discrepancy between observed flow and camera-induced flow. Our objective
keeps this per-pair supervision and adds a multi-frame consistency term
with the \emph{same} static/dynamic split, now defined across a view
graph rather than a single edge.

\paragraph{The view graph and global alignment.}
Collect the views as vertices of a view graph $G=(V,E)$, with an edge
$e=(i,j)$ wherever the backbone relates views $i$ and $j$. Standard fusion
minimizes a confidence-weighted pairwise disagreement
\begin{equation}
  \mathcal{L}_{\mathrm{align}}
  \;=\;
  \sum_{e=(i,j)\in E}\;\sum_{n\in\{i,j\}}\;\sum_{p}
  \nu_n(p)\,
  \big\lVert
     \chi_n(p) \;-\; g_e\!\left(\sigma_e\, X^{(e)}_n(p)\right)
  \big\rVert
  \label{eq:align}
\end{equation}
over per-view $\{D_n,g_n,K_n\}$ and per-edge $\{g_e,\sigma_e\}$, by
gradient descent at inference time~\citep{wang2024dust3r}. Two features
matter. First, \eqref{eq:align} couples variables \emph{only} through
shared correspondence on edges: its structure is a function on the graph.
Second, it is run \emph{post hoc}, per scene, and presumes a consistent
rigid geometry exists---which under object motion it does not. We keep the
edge-local structure of \eqref{eq:align} but move it into training and
replace ``descend on the full residual'' with ``drive the consistent part
to zero and leave the motion part free.''

% --------------------------------------------------------------------
\subsection{Cellular sheaves and the Hodge split}
\label{sec:prelim-sheaves}
% --------------------------------------------------------------------

We recall the minimal vocabulary; see \citet{curry2014sheaves} and
\citet{hansen2019toward} for full treatments.

\begin{definition}[Cellular sheaf]
A \emph{cellular sheaf} $\sheaf$ of finite-dimensional real vector spaces
on $G=(V,E)$ assigns a \emph{stalk} $\sheaf(v)$ to each vertex, a stalk
$\sheaf(e)$ to each edge, and a linear \emph{restriction map}
$\restr{v}{e}:\sheaf(v)\to\sheaf(e)$ to each incident pair
$v\trianglelefteq e$.
\end{definition}

The $0$-cochains $\Cochain^0=\bigoplus_v\sheaf(v)$ collect local data at
every vertex; the $1$-cochains $\Cochain^1=\bigoplus_e\sheaf(e)$ collect
data on every edge. With an orientation on each edge $e=(i,j)$, the
\emph{coboundary} $\coboundary:\Cochain^0\to\Cochain^1$ compares endpoints,
\begin{equation}
  (\coboundary x)_e \;=\; \restr{j}{e}\,x_j \;-\; \restr{i}{e}\,x_i,
  \qquad e=(i,j).
  \label{eq:coboundary}
\end{equation}
A $0$-cochain is a \emph{global section} if $\coboundary x=0$: the local
data agree on every edge after restriction. Equipping stalks with inner
products gives $\coboundary$ an adjoint $\coboundary^{\!\top}$ and the
\emph{sheaf Laplacian}
$\Lap_{\sheaf}=\coboundary^{\!\top}\coboundary$, a sparse PSD operator
inheriting the sparsity of $G$.

\begin{lemma}[Hodge split; obstruction is cohomology
{\citep{hansen2019toward}}]
\label{lem:hodge}
$\ker\Lap_{\sheaf}=\ker\coboundary=\Sections(G;\sheaf)=\Hcoh^0(G;\sheaf)$,
and any $1$-cochain $r\in\Cochain^1$ decomposes orthogonally as
\[
  r \;=\; \underbrace{\coboundary\xi}_{\text{coboundary part}}
        \;\oplus\;
        \underbrace{h}_{\text{harmonic part}},
  \qquad
  h\in\ker\coboundary^{\!\top}\cong\Hcoh^1(G;\sheaf),
\]
where $\coboundary\xi$ is the part of $r$ that some choice of vertex data
(a per-view correction) removes, and $h$ is the irreducible residual no
such choice can remove. In particular $\Hcoh^1=0$ iff every edge residual
is explained by a global section.
\end{lemma}

Lemma~\ref{lem:hodge} is the backbone of the method. Reading the edge
residual of \eqref{eq:align} as a $1$-cochain, the coboundary part is
exactly what a per-view pose/frame correction absorbs---including camera
drift---and the harmonic part is what remains, the cross-view
disagreement attributable to object motion.

% --------------------------------------------------------------------
\subsection{The reconstruction sheaf}
\label{sec:recon-sheaf}
% --------------------------------------------------------------------

Place the fusion problem of \eqref{eq:align} on the sheaf of
Section~\ref{sec:prelim-sheaves}, on the view graph
$G=(V,E)$.\footnote{We instantiate stalks in two ways. A
\emph{coordinate} model takes $\sheaf(v)$ to be the per-pixel world-frame
points $\chi_v$ of \eqref{eq:chi}; a \emph{pose} model takes $\sheaf(v)$
to be the tangent space of view $v$'s frame and scale, a copy of
$\mathfrak{se}(3)\oplus\R$, used where a per-view correction must be
linear. \todo{Decide which model the loss uses in practice; the pose
model makes the coboundary correction exactly ``a camera fix,'' the
coordinate model is simpler to differentiate.}}
The edge stalk $\sheaf(e)$ over $e=(i,j)$ is the geometry on the shared
overlap, into which both endpoints restrict via back-projection
\eqref{eq:chi} and the edge's correspondence field. By construction,
$(\coboundary x)_e=0$ encodes exactly the per-edge agreement term inside
\eqref{eq:align}, so $\mathcal{L}_{\mathrm{align}}$ is, up to confidence
weights folded into the stalk inner products, the sheaf Dirichlet energy
$\tfrac12\inner{x}{\Lap_{\sheaf}x}=\tfrac12\norm{\coboundary x}^2$.

\begin{remark}[Two subtleties carried into training]
\label{rem:subtleties}
First, poses live in a non-abelian group, so the linear coboundary
\eqref{eq:coboundary} is exact only on tangent stalks; in training we
recompute the residual each step, a re-linearization that is automatic
under gradient descent~\citep{singer2011angular,rosen2019se}. Second, the
restriction maps depend on the correspondence field, which depends on the
geometry being solved for; we treat the backbone's per-pair predictions as
the source of restriction maps within a step and let them co-evolve with
the weights across steps.
\end{remark}

% ====================================================================
\section{Method}
\label{sec:method}
% ====================================================================

\paragraph{Overview.}
We fine-tune a $\mathrm{D}^2$USt3R backbone on $N$-frame tuples with an
objective that adds a multi-frame consistency term to the existing
per-pair supervision. The consistency term uses the Hodge split of
Lemma~\ref{lem:hodge}: it drives the coboundary part of the cross-view
residual to zero on static regions, while leaving the harmonic part free
on dynamic regions. A static-scene anchor and a scale gauge make it
non-degenerate. The total objective is
\begin{equation}
  \mathcal{L}
  \;=\;
  \underbrace{\Lstatic + \Ldyn}_{\text{per-pair (D$^2$USt3R)}}
  \;+\;
  \lambda_{\mathrm{c}}\,\Lcons
  \;+\;
  \lambda_{\mathrm{a}}\,\Lanchor,
  \label{eq:total}
\end{equation}
where $\Lstatic,\Ldyn$ are retained from $\mathrm{D}^2$USt3R and the new
terms are defined below. \todo{Confirm the per-pair terms are kept as-is
or chained across the tuple; see E2.}

\paragraph{Multi-frame consistency (the new term).}
For a tuple with view graph $G=(V,E)$, form the per-edge world-frame
residual from the backbone's predictions,
$r_e = \restr{j}{e}x_j - \restr{i}{e}x_i$, and take its harmonic part
$h_e$ via the projection of Lemma~\ref{lem:hodge}
(equivalently, $h = r - \coboundary\xi^\star$ with
$\xi^\star=\arg\min_\xi\norm{r-\coboundary\xi}^2$, a sparse least-squares
solve, differentiable in the predictions). The consistency loss penalizes
the \emph{coboundary} part on static pixels and the \emph{harmonic} part
only where it should not be---i.e. on static pixels---while leaving
dynamic pixels free:
\begin{equation}
  \Lcons
  \;=\;
  \sum_{e\in E}\sum_{p}
  \big(1-\Mdyn(p)\big)\,
  \nu(p)\,
  \norm{\,r_e(p)\,}^2 .
  \label{eq:cons}
\end{equation}
On static pixels the full residual must be coboundary (a camera fix
explains it); on dynamic pixels the gate $\big(1-\Mdyn\big)$ removes the
penalty, so the harmonic mass that motion produces is permitted rather
than fought. \todo{Decide whether dynamic pixels are merely \emph{ungated}
(left free, as written) or \emph{positively supervised} by chaining
$\Ldyn$ across the tuple via flow---the latter is stronger but needs
multi-frame correspondence chaining.}

\paragraph{Non-degeneracy: static anchor + scale gauge.}
On rigid sequences, $\Mdyn\equiv 0$, so \eqref{eq:cons} forces the
\emph{entire} residual---and hence the harmonic mass---toward zero: a
scene with no motion may carry no irreducible inconsistency. This is the
anchor that forbids the network from explaining static structure as
motion. We additionally pin the uniform log-scale direction (the DC mode
of $\Sim(3)$, which otherwise admits a collapse-to-a-point that drives
\eqref{eq:cons} to zero trivially):
\begin{equation}
  \Lanchor
  \;=\;
  \underbrace{\mathbb{1}[\text{static tuple}]\sum_{e,p}\norm{h_e(p)}^2}_{\text{harmonic must vanish}}
  \;+\;
  \underbrace{\big(\log\textstyle\operatorname{med}_{n,p}D_n(p)-\log s_0\big)^2}_{\text{scale gauge}} .
  \label{eq:anchor}
\end{equation}
\todo{The two pieces of \eqref{eq:anchor} may want separate weights; the
scale term may instead be a hard normalization per step.}

\begin{remark}[Why the split is the contribution, not the consistency]
Dropping the gate and the split in \eqref{eq:cons} recovers a naive
multi-frame consistency loss, which penalizes moving-object disagreement
and reproduces the MonST3R smearing failure. The ablation E2 is designed
to show that the split, not multi-frame supervision per se, is what
buys the dynamic-region gain.
\end{remark}

% ====================================================================
\section{Experiments (skeleton)}
\label{sec:experiments}
% ====================================================================

This section is an outline of the claims to establish and the experiments
that test them. Each subsection states the claim, the comparison, the
metric, and---explicitly---the result that would \emph{falsify} the claim,
so a null result is informative rather than ambiguous.

% --------------------------------------------------------------------
\subsection{Setup}
\label{sec:exp-setup}
% --------------------------------------------------------------------

\paragraph{Backbone and training.}
We start from a $\mathrm{D}^2$USt3R checkpoint, freeze the encoder, and
fine-tune the decoder and heads (the protocol $\mathrm{D}^2$USt3R and
MonST3R use), now on $N$-frame tuples rather than pairs.
\todo{Fix $N$ (e.g. $N\!\in\!\{4,5\}$), the view-graph topology
(temporal chain $+$ loop edges), optimizer, LR, epochs, GPU budget.}

\paragraph{Training data and the role of each set.}
\begin{itemize}[leftmargin=1.4em,itemsep=0.1em,topsep=0.1em]
  \item \textbf{PointOdyssey}~\todo{cite}: long realistic-dynamic
        sequences with point tracks; the primary source of $N>2$ dynamic
        supervision and of chained correspondence for the dynamic term.
  \item \textbf{TartanAir}~\todo{cite}: large, fully \emph{static}; the
        non-degeneracy anchor of \eqref{eq:anchor}---harmonic mass must
        vanish here. Used for training the anchor, \emph{not} as a
        benchmark (the backbone has seen it).
  \item \textbf{(Optional) BlinkVision / Spring}~\todo{cite}: realistic
        dynamics with dynamic masks and flow, if additional dynamic
        coverage is needed.
\end{itemize}

\paragraph{Evaluation data (held out).}
\begin{itemize}[leftmargin=1.4em,itemsep=0.1em,topsep=0.1em]
  \item \textbf{Bonn}, \textbf{TUM-Dynamics}~\todo{cite}: real,
        static-dominant dynamic scenes with accurate backbone
        predictions---the clean regime for multi-frame depth.
  \item \textbf{KITTI}~\todo{cite}: driving; static-dominant background
        with rigid moving vehicles---the motivating regime, and the one
        where multi-frame static consistency should help most.
  \item \textbf{Sintel}~\todo{cite}: reported \emph{only} for
        comparability with prior work, with the caveat that its
        dynamic-dominant scenes are out of the regime the method targets;
        not used to support the main claim.
\end{itemize}
\todo{Confirm train/test disjointness for every eval set against the
backbone's training data.}

\paragraph{Metrics.}
Affine-invariant depth (AbsRel, $\delta_1$) over the whole scene and over
dynamic regions separately, following $\mathrm{D}^2$USt3R; camera pose
(rotation/translation error) where applicable.

% --------------------------------------------------------------------
\subsection{E1 --- Headline: multi-frame consistency improves reconstruction}
\label{sec:e1}
% --------------------------------------------------------------------
\textbf{Claim.} Adding $\Lcons$ improves multi-frame reconstruction over
the two-frame $\mathrm{D}^2$USt3R baseline, especially as the number of
frames grows.\\
\textbf{Comparison.} $\mathrm{D}^2$USt3R (two-frame $+$ post-hoc global
alignment) vs.\ ours (trained with $\Lcons$) on $N$-frame tuples.\\
\textbf{Metric.} Multi-frame depth AbsRel/$\delta_1$, all and dynamic.\\
\textbf{Falsifier.} No improvement on \emph{static} regions of held-out
multi-frame sequences (Bonn/TUM/KITTI) $\Rightarrow$ the consistency term
does nothing the post-hoc optimizer was not already doing.

\begin{table}[h]
\centering
\begin{tabular}{lcccc}
\toprule
 & \multicolumn{2}{c}{All} & \multicolumn{2}{c}{Dynamic} \\
Method & AbsRel$\downarrow$ & $\delta_1\uparrow$ & AbsRel$\downarrow$ & $\delta_1\uparrow$ \\
\midrule
$\mathrm{D}^2$USt3R (baseline) & \todo{} & \todo{} & \todo{} & \todo{} \\
Ours ($+\,\Lcons$)            & \todo{} & \todo{} & \todo{} & \todo{} \\
\bottomrule
\end{tabular}
\caption{E1: multi-frame depth on held-out static-dominant sequences.
\todo{Mean over Bonn/TUM/KITTI; report per-dataset in appendix.}}
\label{tab:e1}
\end{table}

% --------------------------------------------------------------------
\subsection{E2 --- The split is what matters (vs.\ naive consistency)}
\label{sec:e2}
% --------------------------------------------------------------------
\textbf{Claim.} The coboundary/harmonic split, not multi-frame
supervision per se, produces the dynamic-region gain; naive multi-frame
consistency \emph{hurts} dynamic regions by smearing moving objects.\\
\textbf{Comparison.} (a) no consistency term (baseline); (b) naive
consistency (penalize the full residual on all pixels, no gate, no split);
(c) ours (gated split, \eqref{eq:cons}).\\
\textbf{Metric.} Dynamic-region depth; qualitative pointmaps on moving
objects.\\
\textbf{Falsifier.} (b) and (c) tie on dynamic regions $\Rightarrow$ the
split is unnecessary and the story collapses to ``multi-frame helps.''

\begin{table}[h]
\centering
\begin{tabular}{lcc}
\toprule
Consistency variant & Dyn.\ AbsRel$\downarrow$ & Dyn.\ $\delta_1\uparrow$ \\
\midrule
(a) none (baseline)        & \todo{} & \todo{} \\
(b) naive (full residual)  & \todo{} & \todo{} \\
(c) split (ours)           & \todo{} & \todo{} \\
\bottomrule
\end{tabular}
\caption{E2: ablating the Hodge split. Expectation: (b) $<$ (a) on dynamic
regions (smearing), (c) $>$ (a).}
\label{tab:e2}
\end{table}

% --------------------------------------------------------------------
\subsection{E3 --- Non-degeneracy: the objective is not gamed}
\label{sec:e3}
% --------------------------------------------------------------------
\textbf{Claim.} The static anchor and scale gauge prevent collapse; on
static-scene controls the network produces near-zero harmonic mass, and
removing the anchors lets the objective degrade.\\
\textbf{Comparison.} ours with vs.\ without (i) the static anchor, (ii)
the scale gauge.\\
\textbf{Metric.} Harmonic energy on held-out static scenes (should be
$\approx 0$); depth scale statistics (no collapse); overall depth.\\
\textbf{Falsifier.} Removing the anchors does \emph{not} degrade results
$\Rightarrow$ the anchors are inert and the gameability concern (and thus
part of the method's justification) was misplaced.
\todo{Also report: does harmonic energy on truly static held-out scenes
stay low without those scenes being in training (generalization of the
anchor)?}

% --------------------------------------------------------------------
\subsection{E4 --- Robustness to frame count and interval}
\label{sec:e4}
% --------------------------------------------------------------------
\textbf{Claim.} The benefit grows (or at least holds) with the number of
frames and with temporal stride, where the post-hoc two-frame baseline
degrades.\\
\textbf{Comparison.} ours vs.\ baseline across $N$ and stride
$\Delta t\in\{1,3,5,7,9\}$ (cf.\ $\mathrm{D}^2$USt3R's interval study).\\
\textbf{Metric.} Multi-frame depth vs.\ $\Delta t$ and $N$.\\
\textbf{Falsifier.} The gap to baseline is flat or shrinking in $N$
$\Rightarrow$ multi-frame structure is not being exploited.

% --------------------------------------------------------------------
\subsection{E5 --- (Optional) Pose-error robustness as a diagnostic}
\label{sec:e5}
% --------------------------------------------------------------------
\textbf{Claim (diagnostic, not headline).} Because camera-pose error lives
in the coboundary, the harmonic part of the residual---and hence the
dynamic supervision derived from it---is robust to pose error, unlike a
raw-residual consistency term.\\
\textbf{Comparison.} harmonic vs.\ raw-residual term under injected camera
drift.\\
\textbf{Metric.} Stability of the term / of dynamic-region depth as drift
grows.\\
\textbf{Note.} This is the property the sheaf formalism predicts and the
one place the cohomological framing is doing visible work; report it as a
mechanism check, not as a competitor to flow.

% --------------------------------------------------------------------
\subsection{What this section is \emph{not} testing}
\label{sec:exp-scope}
% --------------------------------------------------------------------
We do not evaluate an unsupervised dynamic-region detector, do not compare
against optical-flow motion segmentation, and do not claim improved
camera-pose estimation as a headline. The contribution is a training-time
multi-frame consistency objective that improves reconstruction on top of
$\mathrm{D}^2$USt3R; the experiments are scoped to that claim.

% ====================================================================
\bibliographystyle{abbrvnat}
\bibliography{references}
% ====================================================================

\end{document}